% ---------------------------------------------------------------------
\section{Modélisation}
\label{sec:modelisation}
% ---------------------------------------------------------------------

Entre le fichier d'origine $x(t)$ et l'extrait capté $y(t)$ s'intercalent un haut-parleur,
une salle, un microphone, un codec et du bruit ambiant. L'énoncé les résume par
\begin{equation}
  y(t) = \alpha\,(h * x)(t - t_0) + b(t),
  \tag{1}\label{eq:modele}
\end{equation}
auquel s'ajoute le codage perceptuel, qui n'est pas linéaire et ne s'écrit pas sous cette
forme. Chaque terme a une cause physique et appelle une invariance précise de l'empreinte
(\cref{tab:modele}).

\begin{itemize}
  \item $\alpha > 0$ est un \textbf{gain} inconnu : volume de lecture, distance à la source,
        sensibilité du micro. Il est constant à l'échelle d'un extrait.
  \item $h$ est la \textbf{réponse impulsionnelle} de la chaîne haut-parleur $\to$ salle
        $\to$ micro. Sa partie courte (haut-parleur, micro, premières réflexions) colore le
        spectre ; sa partie longue (réverbération tardive) étale le signal dans le temps.
  \item $t_0$ est le \textbf{décalage} inconnu : on ignore où l'extrait a été capté.
  \item $b(t)$ est un \textbf{bruit additif} indépendant de $x$ : souffle blanc, brouhaha
        coloré et non stationnaire.
  \item Le \textbf{codage MP3} supprime les composantes masquées et ajoute un bruit de
        quantification mis en forme sous le seuil de masquage.
\end{itemize}

\begin{table}[h]
  \centering\small
  \caption{Terme du modèle~\eqref{eq:modele}, dégradation qui le simule et invariance visée.}
  \label{tab:modele}
  \begin{tabularx}{\linewidth}{@{}lXX@{}}
    \toprule
    Terme & Dégradation simulée & Invariance visée (où) \\
    \midrule
    $t_0$      & début de requête tiré au hasard, non aligné sur la grille
               & translation du spectrogramme, vote sur $\delta$ (\cref{sec:proprietes,sec:decision}) \\
    $\alpha$   & implicite dans toutes les conditions
               & positions des maxima, seuil local (\cref{sec:maxima}) \\
    $h$ court  & passe-bande \qtyrange{300}{3400}{\hertz}, haut-parleur de téléphone
               & biais $g(k)$ indépendant du temps, condition (5) (\cref{sec:maxima}) \\
    $h$ long   & réverbération $T_{60} = \qty{0.3}{\second}$ et \qty{1}{\second}
               & hors modèle : limite de la propriété (2) (\cref{sec:limites}) \\
    $b$        & bruit blanc (\qtyrange{20}{0}{\decibel}), brouhaha (\qtyrange{10}{0}{\decibel})
               & survie des pics, RSB local par bande (\cref{sec:q6}) \\
    MP3        & \num{128}, \num{64} et \qty{32}{\kilo\bit\per\second}
               & les pics sont les masqueurs (\cref{sec:q6}) \\
    tout       & enregistrements au micro en salle, à deux distances
               & toutes à la fois (\cref{sec:reel}) \\
    \bottomrule
  \end{tabularx}
\end{table}

Les deux exigences sont antagonistes : une empreinte riche (beaucoup de pics, clés fines)
est spécifique mais fragile, une empreinte pauvre est robuste mais confond les morceaux.
Toute la conception (\cref{sec:conception}) consiste à quantifier ce compromis : la
robustesse par la probabilité $p$ qu'un pic survive, la spécificité par le nombre de clés
possibles $M$ et par le nombre d'appariements fortuits qu'il entraîne.
